\section{AHSME 1982}
        \begin{enumerate}
        	\item (\textit{AHSME 1982}) When the polynomial $x^3-2$ is divided by the polynomial $x^2-2$, the remainder is 


 $\textbf{(A)} \ 2 \qquad  \textbf{(B)} \ -2 \qquad  \textbf{(C)} \ -2x-2 \qquad  \textbf{(D)} \ 2x+2 \qquad  \textbf{(E)} \ 2x-2$

 \textbf{Answer: }E



\textbf{Solution}

Working out $\frac{x^3-2}{x^2-2}$ using polynomial long division, we get $x + \frac{2x-2}{x^2-2}$. Thus the answer is $2x -2$, for choice $\boxed{(E)}$.

 



\textbf{Solution 2}

Substituting an easy value, like $3$, you get $\frac{25}{7}$, with a remainder of 4. Then you just plug in x to get $E$

 -dragoon

 


	\item (\textit{AHSME 1982}) If a number eight times as large as $x$ is increased by two, then one fourth of the result equals 


 $\textbf{(A)} \ 2x + \frac{1}{2} \qquad  \textbf{(B)} \ x + \frac{1}{2} \qquad  \textbf{(C)} \ 2x+2 \qquad  \textbf{(D)}\ 2x+4 \qquad \textbf{(E)}\ 2x+16$

 \textbf{Answer: }A



\textbf{Solution}

The number $8$ times as large as $x$ will be $8x$, and $8x$ increased by two will give $8x+2$. Hence finally, the answer is $\frac{1}{4}(8x+2) = \boxed{\text{(A)}\ 2x + \frac{1}{2}}$.

 


	\item (\textit{AHSME 1982}) Evaluate $(x^x)^{(x^x)}$ at $x = 2$. 


 $\textbf{(A)} \ 16 \qquad  \textbf{(B)} \ 64 \qquad  \textbf{(C)} \ 256 \qquad  \textbf{(D)} \ 1024 \qquad  \textbf{(E)} \ 65,536$

 \textbf{Answer: }C



\textbf{Solution}

Plugging in $2$ as $x$ gives $4^4$, which is merely $\boxed{\text{(C)} \ 256 }$.

 


	\item (\textit{AHSME 1982}) The perimeter of a semicircular region, measured in centimeters, is numerically equal to its area, 
measured in square centimeters. The radius of the semicircle, measured in centimeters, is 


 $\textbf{(A)} \ \pi \qquad  \textbf{(B)} \ \frac{2}{\pi} \qquad  \textbf{(C)} \ 1 \qquad  \textbf{(D)} \ \frac{1}{2}\qquad \textbf{(E)} \ \frac{4}{\pi}+2$

 \textbf{Answer: }E



\textbf{Solution}

Let $r$ be the radius of the semicircle.  Then its perimeter is $2r + \pi r = r(2 + \pi)$, while its area is $\frac{1}{2}\pi r^2$.  So:

 $r(2 + \pi) = \frac{1}{2}\pi r^2$
$2 + \pi = \frac{1}{2}\pi r$
$r = \frac{2}{\pi}(2 + \pi) = \boxed{(\mathbf{E})\ \frac{4}{\pi} + 2}$

 -j314andrews

 


	\item (\textit{AHSME 1982}) Two positive numbers $x$ and $y$ are in the ratio $a: b$ where $0 < a < b$. If $x+y = c$, then the smaller of $x$ and $y$ is 


 $\textbf{(A)} \ \frac{ac}{b} \qquad  \textbf{(B)} \ \frac{bc-ac}{b} \qquad  \textbf{(C)} \ \frac{ac}{a+b} \qquad  \textbf{(D)}\ \frac{bc}{a+b}\qquad \textbf{(E)}\ \frac{ac}{b-a}$

 \textbf{Answer: }C



\textbf{Solution}

We can write 2 equations.

 $\frac{x}{y}=\frac{a}{b}$

 and

 $x+y=c$

 Solving for $x$ and $y$ in terms of $a, b, c$ we get :

 $x=\frac{ac}{a+b}$ and $y=\frac{bc}{a+b}$

 Since we know $a$ is less than $b$ and $\frac{x}{y}=\frac{bc}{a+b}$, the smaller of $x$ and $y$ must be $x$. Therefore the answer is $\boxed{\textbf{(C) }\frac{ac}{a+b}}$.

 


 


	\item (\textit{AHSME 1982}) The sum of all but one of the interior angles of a convex polygon equals $2570^\circ$. The remaining angle is 


 $\textbf{(A)} \ 90^\circ \qquad  \textbf{(B)} \ 105^\circ \qquad  \textbf{(C)} \ 120^\circ \qquad  \textbf{(D)}\ 130^\circ\qquad \textbf{(E)}\ 144^\circ$

 \textbf{Answer: }D



\textbf{Solution}

Note that the sum of the interior angles of a convex polygon of $n$ sides is $(n-2) \cdot 180^\circ$, and each interior angle is less than $180^\circ$. Therefore, $n - 2 = \left\lfloor \frac{2570^\circ}{180^\circ} \right\rfloor = 15$, and the missing angle is $15 \cdot 180^\circ - 2570^\circ = 130^\circ \boxed{\text{(D)}}$ .

 


	\item (\textit{AHSME 1982}) If the operation $x \star y$ is defined by $x \star y = (x+1)(y+1) - 1$, then which one of the following is FALSE? 


 $\textbf{(A)} \ x \star y = y\star x  \text{ for all real } x,y. \\ \textbf{(B)} \ x \star (y + z) = ( x \star y ) + (x \star z)  \text{ for all real } x,y, \text{ and } z.\\ \textbf{(C)} \ (x-1) \star (x+1) = (x \star x) - 1 \text{ for all real } x. \\ \textbf{(D)} \ x \star 0 = x \text{ for all real } x. \\ \textbf{(E)} \ x \star (y \star z) = (x \star y) \star z \text{ for all real } x,y, \text{ and } z.$

 \textbf{Answer: }B



\textbf{Solution}

(A) is true because multiplication is commutative.

 (B) is false because we have $x\star (y + z) = (x+1)(y+z+1) - 1$ and $x\star y + x\star z = (x+1)(y+1) - 1 + (x+1)(z+1) - 1$$(x+1)(y+z+2) - 2$$(x+1)(y+z+1) - (x+1) - 2$$[(x+1)(y+z+1) - 1] + x,$which does not match with the previous expression.

 (C) is true because $(x-1)\star (x+1) = ((x-1)+1)((x+1)+1) - 1$$x(x+2) - 1$$x^2 + 2x - 1$$[(x+1)^2 - 1] - 1$$(x\star x) - 1$for all $x$.

 (D) is true because $x\star 0 = (x+1)(0+1) - 1 = (x+1) - 1 = x$ for all $x$.

 (E) is true because multiplication is associative (and the plus-ones inside the parentheses and the minus-ones outside cancel out).

 Therefore, our answer is $\boxed{\textbf{(B)}}$, and we are done.

 


	\item (\textit{AHSME 1982}) By definition, $r! = r(r - 1) \cdots 1$ and $\binom{j}{k} = \frac {j!}{k!(j - k)!}$, where $r,j,k$ are positive integers and $k < j$. 
If $\binom{n}{1}, \binom{n}{2}, \binom{n}{3}$ form an arithmetic progression with $n > 3$, then $n$ equals 


 $\textbf{(A)}\ 5\qquad  \textbf{(B)}\ 7\qquad  \textbf{(C)}\ 9\qquad  \textbf{(D)}\ 11\qquad  \textbf{(E)}\ 12$

 \textbf{Answer: }B



\textbf{Solution}

Since $\binom{n}{1}$, $\binom{n}{2}$, and $\binom{n}{3}$ form an arithmetic progression, $\binom{n}{2} - \binom{n}{1} = \binom{n}{3} - \binom{n}{2}$.  Therefore, 
 \[\frac{n!}{2!\cdot(n-2)!}-\frac{n!}{1!\cdot(n-1)!}=\frac{n!}{3!\cdot(n-3)!}-\frac{n!}{2!\cdot(n-2)!}.\]

 Simplifying these expressions yields $\frac{n(n-1)}{2}-n = \frac {n(n-1)(n-2)}{6}-\frac{n(n-1)}{2}$. Multiplying both sides by $6$ and collecting all terms on one side yields $n^3 - 9n^2 + 14n = 0$, which factors to $n(n-7)(n-2)=0$.  The solutions to this equation are $n \in \{0, 2, 7\}$, but $n > 3$ so the only valid answer is $n = \boxed{(\mathbf{B})\ 7}$.

 


	\item (\textit{AHSME 1982}) A vertical line divides the triangle with vertices $(0,0), (1,1)$, and $(9,1)$ in the $xy\text{-plane}$ into two regions of equal area. 
The equation of the line is $x=$
$\textbf {(A)}\ 2.5 \qquad  \textbf {(B)}\ 3.0 \qquad  \textbf {(C)}\ 3.5 \qquad  \textbf {(D)}\ 4.0\qquad  \textbf {(E)}\ 4.5$

 \textbf{Answer: }B



\textbf{Solution}


                    \begin{figure}[H]
                        \centering
                        \includegraphics[width=0.5\textwidth]{Solution/AHSME_I_1982_9_Solution_Figure_1.png}
                    \end{figure}
                    The vertical line that divides $\triangle ABC$ into two equal regions has equation $x=a$, as shown in the diagram.The area of $ABC$ is $\frac{1}{2}\cdot AC\cdot AF = \frac{1}{2} \cdot 8 \cdot 1 = 4$, so the two regions must each have area $2$ .

 Since $\triangle ABF$ has area $\frac{1}{2}$, the portion of $\triangle ABC$ to the left of $AF$ will be less than $\frac{1}{2}$ and therefore less than $2$. So $\overline{AF}$ is to the left of vertical line $x=a$ (Passing through point $E$).

 The equation of line BC is $y=\frac{1}{9}x$ and the vertical line $x=a$ intersects $\overline{BC}$ at the point $\left(a, \frac{a}{9}\right)$ .  Because the area of the portion of $\triangle ABC$ on the right is $2$, we have 
 \[2=\frac{1}{2}\left(1-\frac{a}{9}\right)(9-a)\] which simplifies to 
 \[(9-a)^2=36\]

 This equation has solutions $a = 3$ and $a = 15$, but $1 < a < 9$ so $a = \boxed{(\mathbf{B}) 3.0}$.

 


	\item (\textit{AHSME 1982}) In the adjoining diagram, $BO$ bisects $\angle CBA$, $CO$ bisects $\angle ACB$, and $MN$ is parallel to $BC$. 
If $AB=12, BC=24$, and $AC=18$, then the perimeter of $\triangle AMN$ is


 
                    \begin{figure}[H]
                        \centering
                        \includegraphics[width=0.5\textwidth]{Problem/AHSME_I_1982_10_Problem_Figure_1.png}
                    \end{figure}
                    
$\textbf {(A)}\ 30 \qquad  \textbf {(B)}\ 33 \qquad  \textbf {(C)}\ 36 \qquad  \textbf {(D)}\ 39 \qquad  \textbf {(E)}\ 42$

 \textbf{Answer: }A



\textbf{Solution}

Since $BO$ and $CO$ are angle bisectors of angles $B$ and $C$ respectively, $\angle MBO = \angle OBC$ and similarly $\angle NCO = \angle OCB$. Because $MN$ and $BC$ are parallel, $\angle OBC = \angle MOB$ and $\angle NOC = \angle OCB$ by corresponding angles. This relation makes $\triangle MOB$ and $\triangle NOC$ isosceles. This makes $MB = MO$ and $NO = NC$. $AM$ + $MB$ = 12, and $AN$ + $NC$ = 18. So, $AM$ + $MO$ = 12, and $AN$ + $NO$ = 18, and those are all of the lengths that make up $\triangle AMN$. Therefore, the perimeter of $\triangle AMN$ is $12 + 18 = 30$. The answer is $\boxed{A}$.

 


	\item (\textit{AHSME 1982}) How many integers with four different digits are there between $1,000$ and $9,999$ such that the absolute value of 
the difference between the first digit and the last digit is $2$? 


 $\textbf {(A)}\ 672 \qquad  \textbf {(B)}\ 784 \qquad  \textbf {(C)}\ 840 \qquad  \textbf {(D)}\ 896 \qquad  \textbf {(E)}\ 1008$

 \textbf{Answer: }C



\textbf{Solution}

All of the digits of the numbers must be 0 to 9. If the first and last digits are x and y, we have $x-y=2$ and $0<x<9$ or $y-x=2,$ and $0<y<9.$ Substituting we have $0<x+2<9,$ and $0<y+2<9.$ Thus $0<x<7$ and $0<y<7,$ which yields 16 pairs (x, y) such that the absolute value of the difference between the x and y is $2.$ However, we are not done. If 0 is the last digit ( with the pair (0, 2):) then we won't have a 4 digit number, so our real value is 15. Because our digits are distinct, there are $(10-2)(10-3)$ ways to fill the middle 2 places with digits, thus by the multiplication principles (counting) there are $15\cdot56 = \boxed {\left(C\right) 840}$ numbers that fulfill these circumstances.

 


	\item (\textit{AHSME 1982}) Let $f(x) = ax^7+bx^3+cx-5$, where $a,b$ and $c$ are constants. If $f(-7) = 7$, then $f(7)$ equals 


 $\textbf {(A)}\ -17 \qquad  \textbf {(B)}\ -7 \qquad \textbf {(C)}\ 14 \qquad  \textbf {(D)}\ 21\qquad  \textbf {(E)}\ \text{not uniquely determined}$

 \textbf{Answer: }A



\textbf{Solution}

$f(x)$ is an odd function shifted down 5 units. Thus, it can be written as $f(x)=g(x)-5$ where $g(x)=ax^7+bx^3+cx$. Thus: $f(-7)=g(-7)-5=7$ and $g(-7)=12$. Using this and the fact $g(x)$ is odd, we can evaluate $f(7)$, which is:

 
 \[f(7) = g(7)-5 = -g(-7)-5 = -12-5 = -17\]

 Therefore, the answer is $\boxed{ \textbf{A}}$.

 


	\item (\textit{AHSME 1982}) If $a>1, b>1$, and $p=\frac{\log_b(\log_ba)}{\log_ba}$, then $a^p$ equals 


 $\textbf {(A)}\ 1 \qquad  \textbf {(B)}\ b \qquad  \textbf {(C)}\ \log_ab \qquad  \textbf {(D)}\ \log_ba \qquad  \textbf {(E)}\ a^{\log_ba}$

 \textbf{Answer: }D



\textbf{Solution 1}

p (log <sub> b</sub> a) = log <sub>b </sub> (log <sub> b</sub> a)

 log <sub> b</sub> (a<sup> p</sup>) =log <sub> b </sub> (log<sub>b</sub> a)

 a<sup>p</sup> = log <sub> b</sub> a

 



\textbf{Solution 2}

Notice that $\frac{\log_b(\log_ba)}{\log_ba}$ strongly resembles the change of base rule. Recall that $\log_ba=\frac{\log_ca}{\log_cb}$. Taking the base on the RHS to be $b$, we get that $p = \log_a(\log_ba)$. Raising $a$ to both sides, we get that 
 \[a^p = a^{\log_a(\log_ba)}\] 
 \[= \boxed{\textbf{(D)} ~ \log_ba}\]



	\item (\textit{AHSME 1982}) In the adjoining figure, points $B$ and $C$ lie on line segment $AD$, and $AB, BC$, and $CD$ are diameters of circle 
$O, N$, and $P$, respectively. Circles $O, N$, and $P$ all have radius $15$ and the line $AG$ is tangent to circle $P$ at $G$. 
If $AG$ intersects circle $N$ at points $E$ and $F$, then chord $EF$ has length


 
                    \begin{figure}[H]
                        \centering
                        \includegraphics[width=0.5\textwidth]{Problem/AHSME_I_1982_14_Problem_Figure_1.png}
                    \end{figure}
                    
$\textbf {(A)}\ 20 \qquad  \textbf {(B)}\ 15\sqrt{2} \qquad  \textbf {(C)}\ 24 \qquad  \textbf {(D)}\ 25 \qquad  \textbf {(E)}\ \text{none of these}$

 \textbf{Answer: }C



\textbf{Solution}

Drop a perpendicular line from $N$ to $AG$ at point $H$. $AN=45$, and since $\triangle{AGP}$ is similar to $\triangle{AHN}$. $NH=9$. $NE=NF=15$ so by the Pythagorean Theorem, $EH=HF=12$. Thus $EF=\boxed{24.}$ Answer is then $\boxed{C}$.

 


	\item (\textit{AHSME 1982}) Let $[z]$ denote the greatest integer not exceeding $z$. Let $x$ and $y$ satisfy the simultaneous equations 


 
 \begin{align*} y&=2[x]+3 \\ y&=3[x-2]+5. \end{align*}



 If $x$ is not an integer, then $x+y$ is 


 $\textbf{(A) } \ \text{ an integer} \qquad  \textbf{(B) } \ \text{ between 4 and 5} \qquad  \textbf{(C) } \ \text{ between  -4 and 4}\qquad\\ \textbf{(D) } \ \text{ between 15 and 16}\qquad \textbf{(E) } \  16.5$

 \textbf{Answer: }D



\textbf{Solution}

We simply ignore the floor of $x$. Then, we have $y$ = $2x + 3$ = $3(x-2)+5$. Solving for $3x - 1 = 2x + 3$, we get $x = 4$. For the floor of $x$, we have $x$ is between $4$ and $5$. Plugging in $8$ + $3$ = $11$ for $y$, we have $y = 11$. We have $11 + 4.x$ = $\boxed {(D)}$




\textbf{Solution 2 (RIGID)}

Since $x$ is not an integer, we let $x=a+b$, where $0<b<1$. 

 So $2[x]+3=2a+3$. $3[x-2]+5=3a-1$.

 $2a+3=3a-1$. $a=4$. So we know that $x$ is between 4 and 5. $y=11$. So $x+y$ is between $15$ and $16$. Select $\boxed{D}$.

 


	\item (\textit{AHSME 1982}) A wooden cube has edges of length $3$ meters. Square holes, of side one meter, centered in each face are cut through to the opposite face. 
The edges of the holes are parallel to the edges of the cube. The entire surface area including the inside, in square meters, is 


 $\textbf {(A)} \ 54 \qquad  \textbf {(B)} \ 72 \qquad  \textbf {(C)} \ 76 \qquad  \textbf {(D)} \ 84\qquad  \textbf {(E)} \ 86$

 \textbf{Answer: }B



\textbf{Solution}

Each exterior unit square which is removed exposes 4 faces of the unit interior squares, so the entire surface area in square meters is $6 \cdot 3^2 - 6 + 24=72.$



	\item (\textit{AHSME 1982}) How many real numbers $x$ satisfy the equation $3^{2x+2}-3^{x+3}-3^{x}+3=0$? 


 $\textbf {(A)}\ 0 \qquad  \textbf {(B)}\ 1 \qquad  \textbf {(C)}\ 2 \qquad  \textbf {(D)}\ 3 \qquad  \textbf {(E)}\ 4$

 \textbf{Answer: }C



\textbf{Solution}

Let $a = 3^x$. Then the preceding equation can be expressed as the quadratic, 
 \[9a^2-28a+3 = 0\] Solving the quadratic yields the roots $3$ and $1/9$. Setting these equal to $3^x$, we can immediately see that there are $\boxed{2}$ real values of $x$ that satisfy the equation.

 


	\item (\textit{AHSME 1982}) In the adjoining figure of a rectangular solid, $\angle DHG=45^\circ$ and $\angle FHB=60^\circ$. Find the cosine of $\angle BHD$.


 
                    \begin{figure}[H]
                        \centering
                        \includegraphics[width=0.5\textwidth]{Problem/AHSME_I_1982_18_Problem_Figure_1.png}
                    \end{figure}
                    
$\textbf{(A)}\ \frac{\sqrt{3}}{6} \qquad  \textbf{(B)}\ \frac{\sqrt{2}}{6} \qquad  \textbf{(C)}\ \frac{\sqrt{6}}{3} \qquad  \textbf{(D)}\ \frac{\sqrt{6}}{4}\qquad \textbf{(E)}\ \frac{\sqrt{6}-\sqrt{2}}{4}$

 \textbf{Answer: }D



\textbf{Solution}

WLOG, let $CD=1$. 

 Looking at square GHDC, we see that $\angle DHC=45$, which implies that $DC=CH=1$ and $DH=\sqrt{2}$ 

 Taking each cross-section one at a time, we look at square DHFB. We obviously know that CHB is a $30$ degree angle, giving $BH=\frac{2\sqrt{3}}{3}$, and $BC=\frac{\sqrt{3}}{3}$.

 Looking at square ABCD, we see that $BD^2=1+\frac{1}{3}\implies BD^2=\frac{4}{3}=\frac{2\sqrt{3}}{3}$. 

 We now look at triangle DBH, which has side lengths $BD=\frac{2\sqrt{3}}{3}$, $BH=\frac{2\sqrt{3}}{3}$ and $DH=\sqrt{2}$. Because $BD$ is opposite angle DHB, by the law of cosines, $\frac{4}{3}=\frac{4}{3}+2-2(\frac{2\sqrt{3}}{3})(\sqrt{2})\cos(\theta)\implies 2=\frac{4\sqrt{6}\cos(\theta)}{3} \implies \frac{2}{\frac{4\sqrt{6}}{3}}\implies \frac{6}{4\sqrt{6}} \implies \frac{\sqrt{6}}{4}$ 

 Thus, the answer is $(\textbf{D})$



	\item (\textit{AHSME 1982}) Let $f(x)=|x-2|+|x-4|-|2x-6|$ for $2 \leq x\leq 8$. The sum of the largest and smallest values of $f(x)$ is 


 $\textbf {(A)}\ 1 \qquad  \textbf {(B)}\ 2 \qquad  \textbf {(C)}\ 4 \qquad  \textbf {(D)}\ 6 \qquad  \textbf {(E)}\ \text{none of these}$

 \textbf{Answer: }B



\textbf{Solution}

Note that at $x=2,3,4,$ one of the three absolute values is equal to $0.$

 Without using absolute values, we rewrite $f(x)$ as a piecewise function:
 \[f(x) = \begin{cases} (x-2)+(4-x)-(6-2x) & \mathrm{if} \ 2\leq x<3 \\ (x-2)+(4-x)-(2x-6) & \mathrm{if} \ 3\leq x<4 \\ (x-2)+(x-4)-(2x-6) & \mathrm{if} \ 4\leq x\leq8 \end{cases},\]which simplifies to
 \[f(x) = \begin{cases} 2x-4 & \mathrm{if} \ 2\leq x<3 \\ -2x+8 & \mathrm{if} \ 3\leq x<4 \\ 0 & \mathrm{if} \ 4\leq x\leq8 \end{cases}.\]The graph of $y=f(x)$ is shown below.
                    \begin{figure}[H]
                        \centering
                        \includegraphics[width=0.5\textwidth]{Solution/AHSME_I_1982_19_Solution_Figure_1.png}
                    \end{figure}
                    The largest value of $f(x)$ is $2,$ and the smallest value of $f(x)$ is $0.$ So, their sum is $\boxed{\textbf {(B)}\ 2}.$



	\item (\textit{AHSME 1982}) The number of pairs of positive integers $(x,y)$ which satisfy the equation $x^2+y^2=x^3$ is 


 $\textbf {(A)}\ 0 \qquad  \textbf {(B)}\ 1 \qquad  \textbf {(C)}\ 2 \qquad  \textbf {(D)}\ \text{not finite} \qquad  \textbf {(E)}\ \text{none of these}$

 \textbf{Answer: }D



\textbf{Solution}

Rearrange the equation to $y^2=x^3-x^2=x^2(x-1)$. This equation is satisfied whenever $x-1$ is a perfect square. There are infinite possible values of $x$, and thus the answer is $\boxed{D: \text{Not Finite}}$



	\item (\textit{AHSME 1982}) In the adjoining figure, the triangle $ABC$ is a right triangle with $\angle BCA=90^\circ$. Median $CM$ is perpendicular to median $BN$, 
and side $BC=s$. The length of $BN$ is


 
                    \begin{figure}[H]
                        \centering
                        \includegraphics[width=0.5\textwidth]{Problem/AHSME_I_1982_21_Problem_Figure_1.png}
                    \end{figure}
                    
$\textbf{(A)}\ s\sqrt 2 \qquad  \textbf{(B)}\ \frac 32s\sqrt2 \qquad  \textbf{(C)}\ 2s\sqrt2 \qquad  \textbf{(D)}\ \frac{s\sqrt5}{2}\qquad \textbf{(E)}\ \frac{s\sqrt6}{2}$

 \textbf{Answer: }E



\textbf{Solution 1}

Suppose that $P$ is the intersection of $\overline{CM}$ and $\overline{BN}.$ Let $BN=x.$ By the properties of centroids, we have $BP=\frac23 x.$

 Note that $\triangle BPC\sim\triangle BCN$ by AA. From the ratio of similitude $\frac{BP}{BC}=\frac{BC}{BN},$ we get
 \begin{align*} BP\cdot BN &= BC^2 \\ \frac23 x\cdot x &= s^2 \\ x^2 &= \frac32 s^2 \\ x &= \boxed{\textbf{(E)}\ \frac{s\sqrt6}{2}}. \end{align*}
~MRENTHUSIASM

 



\textbf{Solution 2}

Let $O$ be the centroid. Knowing that there are two medians, we need to find the length of the third one. Therefore, we draw the median $AP$ such that $O$ is on $AP$. Then, it follows that $BP = CP = OP = \frac{s}{2}$ by Thales's Theorem, and that $AO = s$. So, $AP = \frac{3}{2}s$, which gives us the idea that $CA = s\sqrt{2}$.

 Since $N$ is the median that cuts $CA$, we find out that $CN = AN = \frac{s\sqrt{2}}{2}$. Finally, using Pythagorean again gives $BN=\boxed{{\textbf{(E)}\ \frac{s\sqrt6}{2}}}$.

 


	\item (\textit{AHSME 1982}) In a narrow alley of width $w$ a ladder of length $a$ is placed with its foot at point $P$ between the walls. 
Resting against one wall at $Q$, the distance $k$ above the ground makes a $45^\circ$ angle with the ground. 
Resting against the other wall at $R$, a distance $h$ above the ground, the ladder makes a $75^\circ$ angle with the ground. The width $w$ is equal to 


 $\textbf{(A)}\ a\qquad \textbf{(B)}\ RQ\qquad \textbf{(C)}\ k\qquad \textbf{(D)}\ \frac{h+k}{2}\qquad \textbf{(E)}\ h$

 \textbf{Answer: }E



\textbf{Solution}


                    \begin{figure}[H]
                        \centering
                        \includegraphics[width=0.5\textwidth]{Solution/AHSME_I_1982_22_Solution_Figure_1.png}
                    \end{figure}
                    

 We know that $m\angle QPL=45^{\circ}$ and $m\angle RPT=75^{\circ}.$ Therefore,  $m\angle QPR=60^{\circ}.$ 
 \[\qquad\]Because the two ladders are the same length, we know that
 \[RP=PQ=a.\]Since $\triangle QPR$ is isosceles with vertex angle $60^{\circ},$ we can conclude that it must be equilateral.
 \[\qquad\]Now,  since $\triangle PTR$ is a right triangle and $m\angle TPR=75^{\circ},$ we can conclude that $m\angle PRT=15^{\circ}.$ Because $\triangle QPR$ is equilateral, we know that $m\angle QRP = 60^{\circ}.$It then follows that 
 \[m\angle QRS= m\angle QRP + m\angle PRT = 60^{\circ} + 15^{\circ} = 75^{\circ}.\] Because of ASA,  $\triangle QRS\cong\triangle RPT.$ From there, it follows that $QS=TR=h.$ Since $QS$ is the width of the alley, the answer is $\boxed{\text{E) }h}.$



	\item (\textit{AHSME 1982}) The lengths of the sides of a triangle are consecutive integers, and the largest angle is twice the smallest angle. 
The cosine of the smallest angle is 


 $\textbf{(A)}\ \frac{3}{4}\qquad \textbf{(B)}\ \frac{7}{10}\qquad \textbf{(C)}\ \frac{2}{3}\qquad \textbf{(D)}\ \frac{9}{14}\qquad \textbf{(E)}\ \text{none of these}$

 \textbf{Answer: }A



\textbf{Solution 1 (Law of Sines and Law of Cosines)}

In $\triangle ABC,$ let $a=n,b=n+1,c=n+2,$ and $\angle A=\theta$ for some positive integer $n.$ We are given that $\angle C=2\theta,$ and we need $\cos\theta.$

 We apply the Law of Cosines to solve for $\cos\theta:$ 
 \[\cos\theta=\frac{b^2+c^2-a^2}{2bc}=\frac{n+5}{2(n+2)}.\]We apply the Law of Sines to $\angle A$ and $\angle C:$ 
 \[\frac{\sin\theta}{a}=\frac{\sin(2\theta)}{c}.\]By the Double-Angle Formula $\sin(2\theta)=2\sin\theta\cos\theta,$ we simplify and rearrange to solve for $\cos\theta:$ 
 \[\cos\theta=\frac{c}{2a}=\frac{n+2}{2n}.\]We equate the expressions for $\cos\theta:$ 
 \[\frac{n+5}{2(n+2)}=\frac{n+2}{2n},\]from which $n=4.$ By substitution, the answer is $\cos\theta=\boxed{\textbf{(A)}\ \frac{3}{4}}.$




\textbf{Solution 2 (Law of Cosines Only)}

This solution uses the same variable definitions as Solution 1 does. Moreover, we conclude that $\cos\theta=\frac{n+5}{2(n+2)}$ from the second paragraph of Solution 1.

 We apply the Law of Cosines to solve for $\cos(2\theta):$ 
 \[\cos(2\theta)=\frac{a^2+b^2-c^2}{2ab}=\frac{n-3}{2n}.\]By the Double-Angle Formula $\cos(2\theta)=2\cos^2\theta-1,$ we set up an equation for $n:$ 
 \[\frac{n-3}{2n}=2\left(\frac{n+5}{2(n+2)}\right)^2-1,\] from which $n=-3,-\frac12,4.$ Recall that $n$ is a positive integer, so $n=4.$ By substitution, the answer is $\cos\theta=\boxed{\textbf{(A)}\ \frac{3}{4}}.$



	\item (\textit{AHSME 1982}) In the adjoining figure, the circle meets the sides of an equilateral triangle at six points. If $AG=2, GF=13, FC=1$, 
and $HJ=7$, then $DE$ equals 


 
                    \begin{figure}[H]
                        \centering
                        \includegraphics[width=0.5\textwidth]{Problem/AHSME_I_1982_24_Problem_Figure_1.png}
                    \end{figure}
                    
$\textbf {(A)}\ 2\sqrt{22} \qquad  \textbf {(B)}\ 7\sqrt{3} \qquad  \textbf {(C)}\ 9 \qquad  \textbf {(D)}\ 10 \qquad  \textbf {(E)}\ 13$

 \textbf{Answer: }A



\textbf{Solution 1}

Let $AH=y, BD=a, DE=x,$ and $EC=b.$ Because $AG = 2, GF = 13, HJ = 7$ and $FC = 1$ and they sum to 16, the length of each side of the equilateral triangle will be 16. By Power of a Point on A, we have $y(y + 7) = 2(2 + 13),$ or expanding and factoring, $0 =y^2 + 7y - 30 = (y - 3)(y + 10).$ Therefore, y must equal 3 as this triangle must have positive side lengths.

 Hence, $BJ=16-(3+7)=6.$

 Power of a point on C yields $b(b + x) = 1(1 + 13) = 14,$ and power of a point on B yields a(a + x) = 6(6 + 7) = 78. Also, we know that because a, b, and x are non-intersecting segments of the side of the triangle, a + b + x = 16. 

 So, we have a system of equations:

 
 \[a^2 + ax = 78,\]
 \[b^2 + bx = 14,\]$a + b + x =16.$

 If we subtract the second equation from the first, it appears to be a difference of squares! And we like differences of squares, because we can factor! Subtracting yields $a^2 - b^2 + (a - b)x = (a-b)(a+b)(a-b)x=(a - b)(a + b + x).$ Because we know that $a+b+x=16,$ inputting this we have that $(a - b)16 = 78 - 14 = 64.$ Therefore $a-b=4.$ Then adding this to the third equation, 2a + x = 20, so $a = 10 - (x/2).$ Substituting this into \#1, we can now solve for x, and we have a difference of squares, or $100-x^2/4=78.$ This yields $x^2/4=22,$ so $x/2=\sqrt {22}$ and $x= \left(A\right)2\sqrt{22}!$

 Further clarifications: If you are unsure why we didn't solve for a and b, it is because it is more efficient to use clever manipulations like the ones we used in this problem. We were looking for sums of values we already know, so we indirectly implied dividing $a^2 - b^2 + (a - b)x$ by $a+b+x$ to get $a-b,$ conveniently, then everything fell into place with substitution.

 


	\item (\textit{AHSME 1982}) The adjacent map is part of a city: the small rectangles are blocks, and the paths in between are streets. 
Each morning, a student walks from intersection $A$ to intersection $B$, always walking along streets shown, 
and always going east or south. For variety, at each intersection where he has a choice, he chooses with 
probability $\frac{1}{2}$ whether to go east or south. Find the probability that through any given morning, he goes through $C$. 


 
                    \begin{figure}[H]
                        \centering
                        \includegraphics[width=0.5\textwidth]{Problem/AHSME_I_1982_25_Problem_Figure_1.png}
                    \end{figure}
                    
$\textbf{(A)}\ \frac{11}{32}\qquad \textbf{(B)}\ \frac{1}{2}\qquad \textbf{(C)}\ \frac{4}{7}\qquad \textbf{(D)}\ \frac{21}{32}\qquad \textbf{(E)}\ \frac{3}{4}$

 \textbf{Answer: }D



\textbf{Solution 1}


                    \begin{figure}[H]
                        \centering
                        \includegraphics[width=0.5\textwidth]{Solution/AHSME_I_1982_25_Solution_Figure_1.png}
                    \end{figure}
                    

 The probability that the student passes through $C$ is the sum from $i=0$ to $3$ of the probabilities that he enters intersection $C_i$ in the adjoining figure and goes east. The number of paths from $A$ to $C_i$ is ${{2+i} \choose 2}$, because each such path has $2$ eastward block segments and they can occur in any order. The probability of taking any one of these paths to $C_i$ and then going east is $(\frac{1}{2})^{3+i}$ because there are $3+i$ intersections along the way (including $A$ and $C_i$) where an independent choice with probability $\frac{1}{2}$ is made. So the answer is 
 \[\sum\limits_{i=0}^3 {{2+i} \choose 2} \left( \frac{1}{2} \right)^{3+i} = \frac{1}{8} + \frac{3}{16} + \frac{6}{32} + \frac{10}{64} = \boxed{\frac{21}{32}}.\]



	\item (\textit{AHSME 1982}) If the base $8$ representation of a perfect square is $ab3c$, where $a\ne 0$, then $c$ equals 


 $\textbf{(A)}\ 0\qquad  \textbf{(B)}\ 1 \qquad  \textbf{(C)}\ 3\qquad  \textbf{(D)}\ 4\qquad  \textbf{(E)}\ \text{not uniquely determined}$

 \textbf{Answer: }B



\textbf{Solution}

A perfect square will be $(8k+r)^2=64k^2+16kr+r^2\equiv r^2\pmod{16}$ where $r=0,1,...,7$.

 Notice that $r^2\equiv 1,4,9,0 \pmod{16}$.

 Now $ab3c$ in base 8 is $a8^3+b8^2+3(8)+c\equiv 8+c\pmod{16}$. It being a perfect square means $8+c\equiv 1,4,9,0 \pmod{16}$. That means that c can only be 1 so the answer is 1 = $\boxed{\textbf{(B)}.}$.

 


	\item (\textit{AHSME 1982}) Suppose $z=a+bi$ is a solution of the polynomial equation 
 \[c_4z^4+ic_3z^3+c_2z^2+ic_1z+c_0=0,\] where $c_0, c_1, c_2, c_3, a,$ and $b$ are real constants and $i^2=-1.$ Which of the following must also be a solution? 


 $\textbf{(A)}\ -a-bi\qquad  \textbf{(B)}\ a-bi\qquad  \textbf{(C)}\ -a+bi\qquad  \textbf{(D)}\ b+ai \qquad  \textbf{(E)}\ \text{none of these}$

 \textbf{Answer: }C



\textbf{Solution}

Let $z=wi,$ so the given polynomial equation becomes
 \begin{align*} c_4(wi)^4+ic_3(wi)^3+c_2(wi)^2+ic_1(wi)+c_0&=0 \\ c_4w^4+c_3w^3-c_2w^2-c_1w+c_0&=0. &&(\bigstar) \end{align*}
Note that $(\bigstar)$ is a polynomial equation in $w$ with real coefficients.

 We are given that $w=\frac{a+bi}{i}=b-ai$ is a solution to $(\bigstar).$ By the Complex Conjugate Root Theorem, we conclude that $w=b+ai$ must also be a solution to $(\bigstar),$ from which $z=(b+ai)i=\boxed{\textbf{(C)}\ -a+bi}$ must also be a solution to the given polynomial equation.

 


	\item (\textit{AHSME 1982}) A set of consecutive positive integers beginning with $1$ is written on a blackboard. One number is erased. The average (arithmetic mean) of the remaining numbers is $35\frac{7}{17}$. What number was erased? 


 $\textbf{(A)}\ 6\qquad  \textbf{(B)}\ 7 \qquad  \textbf{(C)}\ 8 \qquad  \textbf{(D)}\  9\qquad  \textbf{(E)}\ \text{cannot be determined}$

 \textbf{Answer: }B



\textbf{Solution}

Suppose that there are $n$ positive integers in the set initially, so their sum is $\frac{n(n+1)}{2}$ by arithmetic series. The average of the remaining numbers is minimized when $n$ is erased, and is maximized when $1$ is erased.

 It is clear that $n>1.$ We write and solve a compound inequality for $n:$
 \begin{alignat*}{8} \frac{\frac{n(n+1)}{2}-n}{n-1} &\leq 35\frac{7}{17} &&\leq \frac{\frac{n(n+1)}{2}-1}{n-1} \\ \frac{n(n+1)-2n}{2(n-1)} &\leq 35\frac{7}{17} &&\leq \frac{n(n+1)-2}{2(n-1)} \\ \frac{n\textasciicircum{}2-n}{2(n-1)} &\leq 35\frac{7}{17} &&\leq \frac{n\textasciicircum{}2+n-2}{2(n-1)} \\ \frac{n(n-1)}{2(n-1)} &\leq 35\frac{7}{17} &&\leq \frac{(n+2)(n-1)}{2(n-1)} \\ \frac{n}{2} &\leq 35\frac{7}{17} &&\leq \frac{n+2}{2} \\ n &\leq 70\frac{14}{17} &&\leq n+2 \\ 68\frac{14}{17} &\leq \hspace{3mm} n &&\leq 70\frac{14}{17}, \end{alignat*}from which $n$ is either $69$ or $70.$

 Let $x$ be the number that is erased. We are given that $\frac{\frac{n(n+1)}{2}-x}{n-1}=35\frac{7}{17},$ or 
 \[\frac{n(n+1)}{2}-x=35\frac{7}{17}\cdot(n-1). \hspace{15mm}(\bigstar)\]

            \begin{itemize}
        \item If $n=70,$ then $(\bigstar)$ becomes $2485-x=\frac{41538}{17},$ from which $x=\frac{707}{17},$ contradicting the precondition that $x$ is a positive integer.
            \end{itemize}
            
            \begin{itemize}
        \item If $n=69,$ then $(\bigstar)$ becomes $2415-x=2408,$ from which $x=\boxed{\textbf{(B)}\ 7}.$
            \end{itemize}
            <u>\textbf{Remark}</u>

 From $(\bigstar),$ note that the left side must be an integer, so must be the right side. It follows that $n-1$ is divisible by $17,$ so $n=69.$



	\item (\textit{AHSME 1982}) Let $x,y,$ and $z$ be three positive real numbers whose sum is $1.$ If no one of these numbers is more than twice any other, 
then the minimum possible value of the product $xyz$ is 


 $\textbf{(A)}\ \frac{1}{32}\qquad \textbf{(B)}\ \frac{1}{36}\qquad \textbf{(C)}\ \frac{4}{125}\qquad \textbf{(D)}\ \frac{1}{127}\qquad \textbf{(E)}\ \text{none of these}$

 \textbf{Answer: }A



\textbf{Solution}

Suppose that the product $xyz$ is minimized at $(x,y,z)=(x_0,y_0,z_0).$ Without the loss of generality, let $x_0 \leq y_0 \leq z_0$ and fix $y=y_0.$

 To minimize $xy_0z,$ we minimize $xz.$ Note that $x+z=1-y_0.$ By a corollary of the AM-GM Inequality <i>\textbf{(If two nonnegative numbers have a constant sum, then their product is minimized when they are as far as possible.)}</i>, we get $z_0=2x_0.$ It follows that $y_0=1-3x_0.$

 Recall that $x_0 \leq 1-3x_0 \leq 2x_0,$ so $\frac15 \leq x_0 \leq \frac14.$ This problem is equivalent to finding the minimum value of 
 \[f(x)=xyz=x(1-3x)(2x)=2x^2(1-3x)\] in the interval $I=\left[\frac15,\frac14\right].$ The graph of $y=f(x)$ is shown below:
                    \begin{figure}[H]
                        \centering
                        \includegraphics[width=0.5\textwidth]{Solution/AHSME_I_1982_29_Solution_Figure_1.png}
                    \end{figure}
                    Since $f$ has a relative minimum at $x=0,$ and cubic functions have at most one relative minimum, we conclude that the minimum value of $f$ in $I$ is at either $x=\frac15$ or $x=\frac14.$ As $f\left(\frac14\right)=\frac{1}{32}\leq f\left(\frac15\right)=\frac{4}{125},$ the minimum value of $f$ in $I$ is $\boxed{\textbf{(A)}\ \frac{1}{32}}.$



	\item (\textit{AHSME 1982}) Find the units digit of the decimal expansion of 
 \[\left(15 + \sqrt{220}\right)^{19} + \left(15 + \sqrt{220}\right)^{82}.\]
$\textbf{(A)}\ 0\qquad  \textbf{(B)}\ 2\qquad  \textbf{(C)}\ 5\qquad  \textbf{(D)}\ 9\qquad  \textbf{(E)}\ \text{none of these}$

 \textbf{Answer: }D



\textbf{Solution 1 (Binomial Expansion)}

Let $A=15+\sqrt{220}$ and $B=15-\sqrt{220}.$ Note that $A^{19}+B^{19}$ and $A^{82}+B^{82}$ are both integers: When we expand (Binomial Theorem) and combine like terms for each expression, the rational terms are added and the irrational terms are canceled.

 We have
 \begin{align*} A^{19}+B^{19} &= \left[\binom{19}{0}15^{19}\sqrt{220}^0+\binom{19}{1}15^{18}\sqrt{220}^1+\cdots+\binom{19}{19}15^0\sqrt{220}^{19}\right] + \left[\binom{19}{0}15^{19}\sqrt{220}^0-\binom{19}{1}15^{18}\sqrt{220}^1+\cdots-\binom{19}{19}15^0\sqrt{220}^{19}\right] \\ &= 2\left[\binom{19}{0}15^{19}\sqrt{220}^0+\binom{19}{2}15^{17}\sqrt{220}^2+\cdots+\binom{19}{18}15^1\sqrt{220}^{18}\right] \\ &= 2\left[\binom{19}{0}15^{19}+\binom{19}{2}15^{17}220+\cdots+\binom{19}{18}15^1 220^9\right]. \end{align*}
Similarly, we have 
 \[A^{82}+B^{82}=2\left[\binom{82}{0}15^{82}+\binom{82}{2}15^{80}220+\cdots+\binom{82}{82}220^{41}\right].\]We add the two equations and take the sum modulo $10:$
 \begin{align*} \left(A^{19}+A^{82}\right)+\left(B^{19}+B^{82}\right) &= 2\Biggl[\binom{19}{0}15^{19}+\phantom{ }\underbrace{\binom{19}{2}15^{17}220+\cdots+\binom{19}{18}15^1 220^9}_{0\pmod{10}}\phantom{ }\Biggr]+2\Biggl[\binom{82}{0}15^{82}+\phantom{ }\underbrace{\binom{82}{2}15^{80}220+\cdots+\binom{82}{82}220^{41}}_{0\pmod{10}}\phantom{ }\Biggr] \\ &\equiv 2\left[\binom{19}{0}15^{19}\right]+2\left[\binom{82}{0}15^{82}\right] \\ &\equiv 2\left[5\right]+2\left[5\right] \\ &\equiv 0\pmod{10}. \end{align*}
It is clear that $0<B^{82}<B^{19}<B<0.5,$ from which $0<B^{19}+B^{82}<1.$ We conclude that the units digit of the decimal expansion of $B^{19}+B^{82}$ is $0.$ Since the units digit of the decimal expansion of $\left(A^{19}+A^{82}\right)+\left(B^{19}+B^{82}\right)$ is $0,$ the units digit of the decimal expansion of $A^{19}+A^{82}$ is $\boxed{\textbf{(D)}\ 9}.$




\textbf{Solution 2 (Characteristic Polynomials)}

Let $a_n = \left(15 + \sqrt{220}\right)^n + \left(15 - \sqrt{220}\right)^n$.  Since $\left(15 + \sqrt{220}\right) + \left(15 - \sqrt{220}\right) = 30$ and $\left(15 + \sqrt{220}\right)\left(15 - \sqrt{220}\right) = 225 - 220 = 5$, by Vieta's rules, $15 + \sqrt{220}$ and $15 - \sqrt{220}$ are roots of the polynomial $x^2 - 30x + 5$.  This must be the characteristic polynomial of $a_n$, and therefore $a_{n} = 30a_{n-1} - 5a_{n-2}$ for all $n \geq 2$. 

 So $a_0 = 2$, $a_1 = 30$, and $a_2 = 30 \cdot 30 - 5 \cdot 2 = 890$.  Since $a_1$ and $a_2$ are divisible by $10$, $a_3$ must also be divisible by $10$.  By similar logic, $a_n$ is divisible by $10$ for every integer $n \geq 1$.  So $a_{19} + a_{82}$ must be divisible by $10$.

 Now $15 - \sqrt{220} < 15 - \sqrt{210.25} = 15 - 14.5 = 0.5$, so $\left(15 - \sqrt{220}\right)^n < (0.5)^n$.  Let $x = \left(15 - \sqrt{220}\right)^{19} + \left(15 - \sqrt{220}\right)^{82}$.  Then $x < (0.5)^{19} + (0.5)^{82} < 0.5 + 0.5 = 1$.

 Therefore, $\left(15 + \sqrt{220}\right)^{19} + \left(15 + \sqrt{220}\right)^{82} = a_{19} + a_{82} - x > a_{19} + a_{82} - 1$, so $\left(15 + \sqrt{220}\right)^{19} + \left(15 + \sqrt{220}\right)^{82}$ must be between $a_{19}+a_{82}-1$ and $a_{19}+a_{82}$.  Since $a_{19} + a_{82}$ is divisible by $10$, $a_{19}+a_{82}-1$ has a units digit of $9$ and $a_{19}+a_{82}$ has a units digit of $0$.  Thus the decimal expansion of $\left(15 + \sqrt{220}\right)^{19} + \left(15 + \sqrt{220}\right)^{82}$ has a units digit of $\boxed{(\mathbf{D})\ 9}$.

 -j314andrews

 


\end{enumerate}